\section*{Geschnetzeltes mit Paprika}
\ihead{}\chead{Personen:4}\ohead{}
\ifoot{Vorbereitungszeit:25 min}\cfoot{}\ofoot{Kochzeit:30 min}
{\Large Zutaten}
\begin{multicols}{2}
\begin{itemize}
\item \SI{500}{g} Schweineschnitzel
\item \SI{2}{EL} Mehl
\item \num{2} Zwiebeln
\item \num{1} grüne Paprikaschote
\item \num{1} rote Paprikaschote
\item \SI{2}{EL} Öl
\item \SI{1}{EL} Paprikapulver, edelsüß
\item \SI{130}{ml} süße Sahne
\item \SI{130}{ml} Brühe
\item \SI{2}{EL} Tomatenmark
\item Salz und Pfeffer nach Belieben
\end{itemize}
% \columnbreak
\end{multicols}
\noindent {\Large Zubereitung}\\
\\
Das Fleisch waschen, trocken tupfen und in ca. \SI{0,5}{cm} dünne Streifen schneiden.
Mehl mit etwas Salz und Pfeffer vermischen und das Fleisch darin wenden.
Die Zwiebeln schälen und in Ringe schneiden.
Die Paprikaschoten entstielen, entkernen, waschen und in Streifen schneiden.
Das Schweinefleisch bei mittlerer Hitze in heißem Öl kräftig von allen Seiten anbraten.
Zwiebelringe und Paprikastreifen zufügen und andünsten.
Paprikapulver darüberstreuen und kurz verrühren.
Sahne und Brühe zufügen und alles verrühren.
Zugedeckt bei schwacher Hitze ca. \SI{20}{min} köcheln lassen.
Das Tomatenmark einrühren und mit Salz und Pfeffer abschmecken.
Das Geschnetzelte mit Reis, Nudeln oder Rösti servieren.
